\subsection{Pradiniai duomenys}
\label{subsec:pradiniai_duomenys}

Pradiniai duomenys susideda iš $699$ įrašų ir $11$ stulpelių. Pirmasis stulpelis yra paciento ID, stulpeliai 2–10 yra požymiai (reikšmės nuo $1$ iki $10$), o paskutinis stulpelis – klasė ($2$ arba $4$). Duomenyse yra $16$ įrašų su trūkstamomis reikšmėmis, pažymėtomis klaustuku „?“. Informaciją apie duomenys galima matyti \ref{tab:pradiniai_duomenys} bei \ref{tab:pradiniu_duomenu_aprasymas} lentelėje.

\begin{table}[H]
    \centering
    \caption{Pradinių duomenų struktūra (pirmieji 9 įrašai bei 1 įrašas su „?“).}
    \begin{tabular}{|c|c|c|c|c|c|c|c|c|c|c|}
        \hline
        ID & $x_1$ & $x_2$ & $x_3$ & $x_4$ & $x_5$ & $x_6$ & $x_7$ & $x_8$ & $x_9$ & Klasė \\
        \hline
        $1000025$ & $5$ & $1$ & $1$ & $1$ & $2$ & $1$ & $3$ & $1$ & $1$ & $2$ \\ \hline
        $1002945$ & $5$ & $4$ & $4$ & $5$ & $7$ & $10$ & $3$ & $2$ & $1$ & $2$ \\ \hline
        $1015425$ & $3$ & $1$ & $1$ & $1$ & $2$ & $2$ & $3$ & $1$ & $1$ & $2$ \\ \hline
        $1016277$ & $6$ & $8$ & $8$ & $1$ & $3$ & $4$ & $3$ & $7$ & $1$ & $2$ \\ \hline
        $1017122$ & $8$ & $10$ & $10$ & $8$ & $7$ & $10$ & $9$ & $7$ & $1$ & $4$ \\ \hline
        $1057013$ & $8$ & $4$ & $5$ & $1$ & $2$ & $?$ & $7$ & $3$ & $1$ & $4$ \\ \hline
        $1041801$ & $5$ & $3$ & $3$ & $3$ & $2$ & $3$ & $4$ & $4$ & $1$ & $4$ \\ \hline
        $1044572$ & $8$ & $7$ & $5$ & $10$ & $7$ & $9$ & $5$ & $5$ & $4$ & $4$ \\ \hline
        $1048672$ & $4$ & $1$ & $1$ & $1$ & $2$ & $1$ & $2$ & $1$ & $1$ & $2$ \\ \hline
        $1096800$ & $6$ & $6$ & $6$ & $9$ & $6$ & $?$ & $7$ & $8$ & $1$ & $2$ \\ \hline
    \end{tabular}
    \label{tab:pradiniai_duomenys}
\end{table}
\begin{table}[H]
    \centering
    \caption{Pradinių duomenų stulpelių aprašymas.}
    \begin{tabular}{|c|l|c|}
        \hline
        Stulpelis & Aprašymas & Reikšmių intervalas \\
        \hline
        ID & Paciento identifikatorius & Teigiamas sveikasis skaičius \\ \hline
        $x_1$ & Clump Thickness & $1–10$ \\ \hline
        $x_2$ & Uniformity of Cell Size & $1–10$ \\ \hline
        $x_3$ & Uniformity of Cell Shape & $1–10$ \\ \hline
        $x_4$ & Marginal Adhesion & $1–10$ \\ \hline
        $x_5$ & Single Epithelial Cell Size & $1–10$ \\ \hline
        $x_6$ & Bare Nuclei & $1–10$ \\ \hline
        $x_7$ & Bland Chromatin & $1–10$ \\ \hline
        $x_8$ & Normal Nucleoli & $1–10$ \\ \hline
        $x_9$ & Mitoses & $1–10$ \\ \hline
        Klasė & $2$ – nepiktybinis, $4$ – piktybinis & $2$ arba $4$ \\ \hline
    \end{tabular}
    \label{tab:pradiniu_duomenu_aprasymas}
\end{table}

Iš viso pradinėje duomenų aibėje yra 458 nepiktybinių (angl.~\textit{Benign}) (klasė $2$) ir 241 piktybinių (angl.~\textit{Malignant}) (klasė $4$) navikų įrašai. Šeštame stulpelyje (Bare Nuclei) yra 16 trūkstamų reikšmių, pažymėtų klaustuku „?“.